\clearpage
\VerdierPhysicalPageAnchor{0175}%
\begingroup\small\setlength{\parskip}{0pt}\setlength{\emergencystretch}{1.5em}
\setlength{\abovedisplayskip}{4pt plus 2pt minus 2pt}%
\setlength{\belowdisplayskip}{4pt plus 2pt minus 2pt}%

\VerdierSemanticAnchor{ch3d-1-1-2-cont}
由 $K(\mathcal A)$ 中特异三角的定义（第一章，3.3.1）以及典范上同调函子
$H\colon K(\mathcal A)\longrightarrow\mathcal A$ 的性质，立即可知范畴
$K^{*,*'}(\mathcal A)\subset K(\mathcal A)$ 是 $K(\mathcal A)$ 的满三角
子范畴（第二章，2.1.5；当 $*={}$“空”时为严格满），并且是饱和子范畴
（第二章，2.1.6）。

\VerdierSemanticAnchor{ch3d-1-1-3}
\medskip\noindent\textbf{1.1.3.}\quad
若对每个整数 $i$ 都有 $H^i(X)=0$，则称 $\mathcal A$ 的复形 $X$
为\emph{无环复形}。以 $\operatorname{Ac}^{*}(\mathcal A)$ 表示
$K^{*}(\mathcal A)$ 的如下满子范畴：其对象为属于
$K^{*}(\mathcal A)$ 的无环复形。若 $K(\mathcal A)$（相应地，
$\operatorname{comp}(\mathcal A)$）中的态射
$X\xrightarrow{u}Y$ 对每个整数 $i$ 所诱导的态射
$H^i(X)\xrightarrow{H^i(u)}H^i(Y)$ 都是同构，则称该态射为一个
\emph{拟同构}。以 $\operatorname{Qis}^{*}(\mathcal A)$（在不会混淆时
也简记为 $\operatorname{Qis}^{*}$）表示 $K^{*}(\mathcal A)$ 中所有
拟同构的集合。$K^{*}(\mathcal A)$ 的子范畴
$\operatorname{Ac}^{*}(\mathcal A)$ 是 $K^{*}(\mathcal A)$ 中严格满、饱和的
三角子范畴。
$\operatorname{Qis}^{*}(\mathcal A)$ 是 $K^{*}(\mathcal A)$ 的态射组成的
饱和乘法系（第二章，2.1.4）。在 $K^{*}(\mathcal A)$ 的特异三角
$(X,Y,Z,u,v,w)$ 中，态射 $u$ 是拟同构，当且仅当复形 $Z$ 无环。
换言之（第二章，2.1.15），$\operatorname{Qis}^{*}(\mathcal A)$ 是
$K^{*}(\mathcal A)$ 中与满三角子范畴
$\operatorname{Ac}^{*}(\mathcal A)$ 相联系的乘法系。

\VerdierSemanticAnchor{ch3d-1-1-4}
\medskip\noindent\textbf{1.1.4.}\quad
设 $M$ 为 $\mathcal A$ 的一族对象。若 $\mathcal A$ 的复形的每个分量都属于
$M$，则称该复形为\emph{$M$ 型复形}。若 $K^{*}(\mathcal A)$ 或
$\operatorname{comp}^{*}(\mathcal A)$ 中的态射
$u\colon X\longrightarrow Y$ 是拟同构，并且 $Y$（相应地，$X$）是
$M$ 型复形，则称 $u$ 为 $X$（相应地，$Y$）的一个\emph{$M$ 型右分解}
（相应地，\emph{左分解}）。通常，在不会引起混淆时，我们滥用语言，省略
“右”（相应地，“左”）字样。同样，通常在不会引起混淆时，我们也滥用语言，
把 $Y$ 的一个右分解（相应地，左分解）的靶（相应地，源）称为 $Y$ 的一个
右分解（相应地，左分解）。以
$\operatorname{comp}_{\mathcal A}^{*}(M)$，或简记为
$\operatorname{comp}^{*}(M)$（相应地，$K_{\mathcal A}^{*}(M)$，或简记为
$K^{*}(M)$），表示 $\operatorname{comp}^{*}(\mathcal A)$（相应地，
$K^{*}(\mathcal A)$）中由 $M$ 型复形定义的满子范畴。

\VerdierSemanticAnchor{ch3d-1-1-5}
\medskip\noindent\textbf{1.1.5.}\quad
当 $M$ 是范畴 $\mathcal A$ 的投射对象（相应地，内射对象）的集合时，分量
都是投射对象（相应地，内射对象）的复形，依照上述约定，称为\emph{投射型
复形}（相应地，\emph{内射型复形}）。读者应注意，不要把投射型复形
（相应地，内射型复形）与范畴 $\operatorname{comp}(\mathcal A)$ 的投射对象
（相应地，内射对象）混淆。容易验证，$\operatorname{comp}(\mathcal A)$ 的
对象是投射对象（相应地，内射对象），当且仅当它是投射型（相应地，内射型）
并且同伦于零。
\endgroup
